% ---
% id: foundations/req:declared-lattice
% kind: req
% title: "The declared lattice condition"
% summary: >-
%   Every entry of operator data declared at a lattice exponent k is a lattice
%   value at k, and so a binary64 number. The condition binds data an instance
%   snaps or declares to a consumer that admits only at its lattice, and the
%   emission of a producer that declares one, and no other data.
% prominence: standard
% status: implemented
% uses: [def:binary64, lem:binary64-lattice, def:declared-data, def:snap]
% mentions: [req:boundary-admission]
% assumes: []
% refs: []
% provenance: newly-derived
% implements: certiqs_key.foundations.snap:is_on_lattice
% lean:
% ---

\begin{requirement}\label{req:declared-lattice}
(D-LAT) Let $k$ be the one \tref{lattice exponent}
(Definition~\ref{def:binary64}\,(iv)) at which a declarer declares operator
data. Every entry of every operator of those data is a \tref{lattice value} at
$k$: the real and the imaginary part of the entry are integer multiples of
$2^{-k}$ and have modulus at most $1$. Each entry is then exactly a binary64
number, by Lemma~\ref{lem:binary64-lattice}\,(i).

Operator data are declared at a lattice exponent $k$, and (D-LAT) binds them,
in these cases and no other: the data an instance snaps at $k$
(Definition~\ref{def:snap}); the data an instance declares to a consumer that
admits \tref{declared data} only at a declared lattice exponent $k$; and every
operator that a producer which declares a lattice exponent $k$ emits
(Definition~\ref{def:declared-data}). Where data of the second case are
numbers of a supplier's output that the instance passes on unchanged
(Definition~\ref{def:declared-data}\,(ii)), the instance does not establish
(D-LAT) of them: the consumer's exact decision on entry establishes it, and
the consumer refuses data that fail it
(Requirement~\ref{req:boundary-admission}). A producer that declares no
lattice exponent has no obligation under (D-LAT).
\end{requirement}
